% Abstract: 150-250 words summarizing your entire paper
% Include: problem, methods, key findings, contribution

Wikipedia aims to be a public encyclopedia, a place where people can share knowledge for the greater good. It works as a global digital public square for universal contributions of knowledge. While its intentions- and those of most contributors– are noble, a few bad actors can move such a resource from being a public square– one where the sharing of information is done respectfully for the pursuit of the truth– to a tool to spread propaganda, misinformation and defamation for their own personal gain. The unique open editing model allows anyone to contribute, but also offers an interesting outlook into trust and anonymity. Wikipedia’s edits are tracked, and the Wikipedia API allows open access into the IP addresses of unregistered users, in part to prevent vandalism. While this helps in terms of moderating Wikipedia, it raises questions of how free speech plays a role, especially for editors who may live in countries under repressive political regimes. This research investigates the relationship between political context, free speech, and participation on Wikipedia. This informs and increases understanding on platform governance, user trust, and privacy; specifically how politics can change participation in digital public squares.


